\documentclass[a4paper,10pt]{article}

\usepackage[utf8]{inputenc}
\usepackage[T1]{fontenc}
\usepackage[brazil]{babel}
\usepackage{geometry}
\usepackage{parskip}
\usepackage{microtype}
\usepackage{enumitem}
\usepackage{titlesec}
\usepackage{array}
\usepackage{tabularx}
\usepackage{hyperref}
\usepackage{bookmark}
\usepackage{xurl}

\geometry{top=1.3cm, bottom=1.3cm, left=1.4cm, right=1.4cm}
\setcounter{secnumdepth}{0}
\setlist[itemize]{
    leftmargin=1.1em,
    itemsep=0.18em,
    topsep=0.2em,
    parsep=0em,
    partopsep=0em
}
\pagestyle{empty}

\hypersetup{
    pdftitle={Breno Juan de Oliveira Pinto - Engenheiro de Dados / Analista de Dados},
    pdfauthor={Breno Juan de Oliveira Pinto},
    colorlinks=true,
    linkcolor=black,
    urlcolor=black,
    citecolor=black,
    bookmarksdepth=2
}

\titleformat{\section}
{\large\bfseries}
{}
{0em}
{}
[\titlerule\vspace{0.4ex}]
\titlespacing*{\section}{0pt}{0.7em}{0.35em}

\newcounter{cventry}
\newcommand{\cventry}[4]{%
    \refstepcounter{cventry}%
    \phantomsection%
    \pdfbookmark[2]{#1}{cventry-\thecventry}%
    \noindent\begin{tabularx}{\textwidth}{@{}>{\raggedright\arraybackslash}X >{\raggedleft\arraybackslash}X@{}}
    \textbf{#1} & #2 \\
    \textit{#3} & \textit{#4} \\
    \end{tabularx}
}

\begin{document}

\begin{center}
    {\LARGE \textbf{Breno Juan de Oliveira Pinto}} \\
    \vspace{0.1cm}
    Engenheiro de Dados \textbar{} Analista de Dados \\
    \vspace{0.1cm}
    Sao Paulo, Brasil
    {\textbullet}
    \href{mailto:brenojuandeoliveirapinto@gmail.com}{brenojuandeoliveirapinto@gmail.com}
    {\textbullet}
    \href{https://www.linkedin.com/in/brenool1veira}{linkedin.com/in/brenool1veira}
    {\textbullet}
    \href{https://portfolio-breno-oliveira.netlify.app}{portfolio-breno-oliveira.netlify.app}
\end{center}

\section{Resumo Profissional}
Profissional de dados com experiencia pratica na construcao de rotinas ETL/ELT, automacao de processos, integracao de sistemas corporativos e estruturacao de dados para relatorios e analises. Atuo com SQL, Python, C\#, APIs REST/SOAP, conceitos de Data Warehouse e Power BI, com experiencia direta em ambientes ERP como TOTVS Protheus e TOTVS RM. Meu foco esta em transformar dados operacionais em pipelines confiaveis, visibilidade de negocio e insumos acionaveis para areas como Financeiro, Contabilidade, Compras, Estoque, RH e servicos compartilhados.

\section{Experiencia}
\cventry{JOFEGE}{Brasil}{Analista de Sistemas}{Abr 2022 -- Atual}
\begin{itemize}
    \item Desenvolvo e mantenho rotinas de ETL/ELT, integracoes e fluxos de dados que apoiam necessidades analiticas e processos operacionais no ecossistema ERP.
    \item Crio solucoes de automacao e scripts de processamento de dados com SQL, Python e C\# para reduzir atividades manuais e aumentar a consistencia das operacoes recorrentes.
    \item Integro sistemas internos e externos por meio de APIs REST, servicos SOAP e web services do TOTVS RM, melhorando a troca de dados entre aplicacoes.
    \item Atendo demandas de dados e processos das areas Financeira, Contabil, Compras, Estoque, RH e Servicos Globais, conectando regras de negocio com implementacao tecnica.
    \item Estruturo e apoio dashboards, relatorios e bases de dados voltados a acompanhamento, tomada de decisao e visibilidade operacional.
\end{itemize}

\section{Competencias Tecnicas}
\begin{itemize}
    \item \textbf{Engenharia de Dados:} ETL/ELT, pipelines de dados, integracao de dados, modelagem de dados, Data Warehouse, Data Lake, automacao de processos
    \item \textbf{Programacao:} SQL, T-SQL, Python, C\#
    \item \textbf{Bancos de Dados:} Oracle, MySQL, NoSQL, consulta e manipulacao de dados relacionais e nao relacionais
    \item \textbf{Analytics \& BI:} Power BI, criacao de dashboards, relatorios, visualizacao de dados
    \item \textbf{Integracao:} APIs REST, SOAP, integracao entre sistemas, web services do TOTVS RM
    \item \textbf{Dominio ERP:} TOTVS Protheus, TOTVS RM, suporte a processos financeiros e backoffice
\end{itemize}

\section{Formacao Academica}
\cventry{Universidade Federal do Rio de Janeiro}{Brasil}{MBA em Engenharia de Dados}{Em andamento}
\cventry{Universidade Sao Francisco}{Brasil}{Pos-graduacao em Ciencia de Dados}{Concluido}
\cventry{Universidade Sao Francisco}{Brasil}{Tecnologo em Analise e Desenvolvimento de Sistemas}{Concluido}

\section{Idiomas}
\begin{itemize}
    \item \textbf{Portugues:} Nativo
    \item \textbf{Ingles:} Avancado
\end{itemize}

\end{document}
